% ==========================================================================
% Kapitel 3: Übungsblätter & Präsenzaufgaben
% ==========================================================================
\clearpage
\section{Übungen}

\subsection{Übungsblatt 3 und Präsenzaufgaben (17.\ November 2026)}
\vorlesungsdatum{17.\ November 2026}{handschrift/Uebungsblatt 3_261117_094748.PDF}
\notizquelleDatei{handschrift/Uebungsblatt 3_261117_094748.PDF}{1--11}

\subsubsection*{Aufgabe 1: Quantoren, Negation und Alltagssprache}
\notizquelleDatei{handschrift/Uebungsblatt 3_261117_094748.PDF}{1--5}

\noindent\textbf{Aufgabenstellung.}
Formulieren Sie die folgenden Sätze mithilfe von Quantoren und geeignet definierten Aussagen und Aussageformen in formale Aussagen um. Benutzen Sie die Äquivalenzen
\[
    \neg\bigl(\forall x\in\Omega : H(x)\bigr) \iff \exists x\in\Omega : \neg H(x),
    \qquad
    \neg\bigl(\exists x\in\Omega : H(x)\bigr) \iff \forall x\in\Omega : \neg H(x),
\]
um die gebildeten formalen Aussagen umzuformen, und übersetzen Sie das Ergebnis wieder in alltägliche Aussagen.

\begin{enumerate}
    \item \emph{„Es gibt Blätter, die nicht auf Bäumen wachsen.“}
    \begin{itemize}
        \item \textbf{Definitionen:} Sei $B$ die Menge aller Blätter und $H(b)$ die Aussageform „$b$ wächst auf einem Baum“.
        \item \textbf{Formalisierung und Umformung:}
        \[
            \exists b\in B : \neg H(b)
            \iff \neg\bigl(\forall b\in B : H(b)\bigr).
        \]
        \item \textbf{Alltägliche Aussage:} „Es ist nicht wahr, dass alle Blätter auf Bäumen wachsen.“
    \end{itemize}

    \item \emph{„Wenn ich Eintopf koche, benutze ich alle Zutaten, die ich in der Küche finde und die noch haltbar sind.“}
    \begin{itemize}
        \item \textbf{Definitionen:} Seien $A$ die Menge aller Zutaten in meiner Küche und $B$ die Menge aller Zutaten, die noch haltbar sind. Setze $\Omega := A\cap B$, und sei $H(x)$ die Aussageform „$x$ ist in meinem selbstgekochten Eintopf“.
        \item \textbf{Formalisierung und Umformung:}
        \[
            \forall x\in\Omega : H(x)
            \iff \neg\bigl(\exists x\in\Omega : \neg H(x)\bigr)
            \iff \nexists x\in\Omega : \neg H(x).
        \]
        \item \textbf{Alltägliche Aussage:} „Es gibt keine Zutat aus meiner Küche, die noch haltbar ist und nicht in meinem Eintopf ist.“
    \end{itemize}

    \item \emph{„Für jeden handgestrickten Pullover gibt es einen Menschen, der den Pullover gestrickt hat.“}
    \begin{itemize}
        \item \textbf{Definitionen:} Sei $\Omega$ die Menge aller handgestrickten Pullover, $\Omega'$ die Menge aller Menschen und $H(x,y)$ die Aussageform „$x$ wurde von $y$ gestrickt“.
        \item \textbf{Formalisierung und Umformung:}
        \[
            \forall x\in\Omega\;\exists y\in\Omega' : H(x,y)
            \iff \neg\bigl(\exists x\in\Omega\;\forall y\in\Omega' : \neg H(x,y)\bigr).
        \]
        \item \textbf{Alltägliche Aussage:} „Es ist nicht wahr, dass es einen handgestrickten Pullover gibt, der von keinem Menschen gestrickt worden ist.“
    \end{itemize}
\end{enumerate}

\subsubsection*{Aufgabe 2: Mengenidentitäten}
\notizquelleDatei{handschrift/Uebungsblatt 3_261117_094748.PDF}{1, 6--7}

\noindent\textbf{Aufgabenstellung.}
Beweisen oder widerlegen Sie für alle Mengen $A$, $B$ und $C$:
\begin{enumerate}
    \item $\bigl((A\cap B)\setminus C\bigr)\cup\bigl((B\cap C)\setminus A\bigr)\cup\bigl((A\cap C)\setminus B\bigr)
    = \bigl((A\cap B)\cup(B\cap C)\cup(A\cap C)\bigr)\setminus(A\cap B\cap C)$.
    \item $\bigl((A\cup B)\setminus C\bigr)\cap\bigl((B\cup C)\setminus A\bigr)\cap\bigl((A\cup C)\setminus B\bigr)
    = \bigl((A\cup B)\cap(B\cup C)\cap(A\cup C)\bigr)\setminus(A\cap B\cap C)$.
\end{enumerate}

\begin{proof}[Lösung zu Teil 1 (wahr)]
Für beliebige Mengen $X,Y$ gilt $(X\cap Y)\setminus Z = (X\cap Y)\setminus(X\cap Y\cap Z)$, denn ein Element $x\in X\cap Y$ liegt genau dann nicht in $Z$, wenn es nicht in $X\cap Y\cap Z$ liegt. Wendet man dies auf die drei Terme der linken Seite an (in einer festen Obermenge $U$, z.\,B.\ $U=A\cup B\cup C$, mit Komplement $(\cdot)^c$), so erhält man:
\begin{align*}
    \bigl((A\cap B)\setminus C\bigr)\cup\bigl((B\cap C)\setminus A\bigr)\cup\bigl((A\cap C)\setminus B\bigr)
    &= \bigl((A\cap B)\cap (A\cap B\cap C)^c\bigr) \\
    &\quad \cup \bigl((B\cap C)\cap (A\cap B\cap C)^c\bigr) \\
    &\quad \cup \bigl((A\cap C)\cap (A\cap B\cap C)^c\bigr) \\
    &= \bigl((A\cap B)\cup(B\cap C)\cup(A\cap C)\bigr)\cap (A\cap B\cap C)^c \\
    &= \bigl((A\cap B)\cup(B\cap C)\cup(A\cap C)\bigr)\setminus(A\cap B\cap C).
\end{align*}
Anschaulich beschreibt dies genau diejenigen Elemente, die in mindestens zwei der Mengen $A,B,C$, aber nicht in allen dreien liegen:

\begin{center}
\begin{tikzpicture}[scale=0.9]
    % A\cap B
    \begin{scope}
        \clip (0,0.65) circle (1.05);
        \fill[primary!25] (-0.65,-0.4) circle (1.05);
    \end{scope}
    % B\cap C
    \begin{scope}
        \clip (-0.65,-0.4) circle (1.05);
        \fill[primary!25] (0.65,-0.4) circle (1.05);
    \end{scope}
    % A\cap C
    \begin{scope}
        \clip (0,0.65) circle (1.05);
        \fill[primary!25] (0.65,-0.4) circle (1.05);
    \end{scope}
    % Ausschluss A\cap B\cap C
    \begin{scope}
        \clip (0,0.65) circle (1.05);
        \clip (-0.65,-0.4) circle (1.05);
        \fill[white] (0.65,-0.4) circle (1.05);
    \end{scope}
    % Kreise zeichnen
    \draw[thick, primary] (0,0.65) circle (1.05);
    \draw[thick, primary] (-0.65,-0.4) circle (1.05);
    \draw[thick, primary] (0.65,-0.4) circle (1.05);
    \node at (0,1.95) {$A$};
    \node at (-1.9,-0.7) {$B$};
    \node at (1.9,-0.7) {$C$};
\end{tikzpicture}
\end{center}
\end{proof}

\begin{proof}[Widerlegung zu Teil 2 (falsch)]
Wähle als Gegenbeispiel $A=\{1\}$, $B=\{1\}$ und $C=\emptyset$.
\begin{itemize}
    \item \textbf{Linke Seite:} Es ist $(A\cup B)\setminus C = \{1\}$, $(B\cup C)\setminus A = \emptyset$ und $(A\cup C)\setminus B = \emptyset$. Damit ergibt sich
    \[
        \{1\}\cap\emptyset\cap\emptyset = \emptyset.
    \]
    \item \textbf{Rechte Seite:} Es ist $A\cup B = B\cup C = A\cup C = \{1\}$ sowie $A\cap B\cap C = \emptyset$. Damit ergibt sich
    \[
        (\{1\}\cap\{1\}\cap\{1\})\setminus\emptyset = \{1\}.
    \]
\end{itemize}
Da $\emptyset\ne\{1\}$, ist die Aussage im Allgemeinen falsch. (Tatsächlich ist die linke Seite für beliebige Mengen $A,B,C$ stets leer, da ein Element im Schnitt weder in $C$, noch in $A$, noch in $B$ liegen dürfte, zugleich aber in $A\cup B$.)
\end{proof}

\subsubsection*{Präsenzaufgabe 1: Bildmengen von Schnitten}
\notizquelleDatei{handschrift/Uebungsblatt 3_261117_094748.PDF}{8--9}

\noindent\textbf{Aufgabenstellung.}
Gegeben sei die Abbildung $f:\R\to\R$, $x\mapsto\sin(x)$, sowie die Teilmengen
\[
    A_1 = \left\{x\in\R \;\middle|\; x\ge -\frac{\pi}{4}\right\},
    \qquad
    A_2 = \left\{x\in\R \;\middle|\; x\le \frac{\pi}{4}\right\}.
\]
Bestimmen Sie $f(A_1)\cap f(A_2)$ und $f(A_1\cap A_2)$, und prüfen Sie, ob Gleichheit gilt.

\medskip
\noindent\textbf{Lösung.}
Da $\sin$ auf jedem Intervall der Länge $2\pi$ alle Werte in $[-1,1]$ annimmt, gilt:
\[
    f(A_1) = \{f(x)\mid x\ge -\tfrac{\pi}{4}\} = [-1,1],
    \qquad
    f(A_2) = \{f(x)\mid x\le \tfrac{\pi}{4}\} = [-1,1],
\]
und somit
\[
    f(A_1)\cap f(A_2) = [-1,1].
\]
Andererseits ist
\[
    A_1\cap A_2 = \left\{x\in\R \;\middle|\; -\frac{\pi}{4}\le x\le\frac{\pi}{4}\right\} = \left[-\frac{\pi}{4},\frac{\pi}{4}\right].
\]
Auf $\left[-\frac{\pi}{4},\frac{\pi}{4}\right]$ ist $\sin$ streng monoton wachsend mit $\sin\left(-\frac{\pi}{4}\right)=-\frac{\sqrt{2}}{2}$ und $\sin\left(\frac{\pi}{4}\right)=\frac{\sqrt{2}}{2}$, also gilt
\[
    f(A_1\cap A_2) = \left[-\frac{\sqrt{2}}{2},\frac{\sqrt{2}}{2}\right].
\]
Folglich ist $f(A_1)\cap f(A_2)\ne f(A_1\cap A_2)$ (es gilt lediglich die echte Inklusion $f(A_1\cap A_2)\subsetneq f(A_1)\cap f(A_2)$).

\begin{center}
\begin{tikzpicture}[xscale=0.95, yscale=1.55]
    % Markierter Bereich A1 cap A2
    \fill[primary!12] (-0.7854,-0.7071) rectangle (0.7854,0.7071);
    % Achsen
    \draw[-{Stealth[length=2.2mm]}] (-2.3,0) -- (5.1,0) node[right] {$x$};
    \draw[-{Stealth[length=2.2mm]}] (0,-1.25) -- (0,1.25) node[above] {$y$};
    % Sinus-Kurve
    \draw[thick, primary, domain=-2.2:4.9, samples=120] plot (\x, {sin(\x r)});
    % Hilfslinien fuer +-pi/4
    \draw[dashed, secondary] (-0.7854,0) -- (-0.7854,-0.7071) -- (0,-0.7071);
    \draw[dashed, secondary] (0.7854,0) -- (0.7854,0.7071) -- (0,0.7071);
    % Ticks x-Achse
    \draw (-1.5708,0.06) -- (-1.5708,-0.06) node[below] {\scriptsize$-\frac{\pi}{2}$};
    \draw (-0.7854,0.06) -- (-0.7854,-0.06) node[above=2pt] {\scriptsize$-\frac{\pi}{4}$};
    \draw (0.7854,0.06) -- (0.7854,-0.06) node[below=2pt] {\scriptsize$\frac{\pi}{4}$};
    \draw (1.5708,0.06) -- (1.5708,-0.06) node[below] {\scriptsize$\frac{\pi}{2}$};
    \draw (4.7124,0.06) -- (4.7124,-0.06) node[below] {\scriptsize$\frac{3\pi}{2}$};
    % Ticks y-Achse
    \draw (-0.06,1) -- (0.06,1) node[left=3pt] {\scriptsize$1$};
    \draw (-0.06,0.7071) -- (0.06,0.7071) node[left=3pt] {\scriptsize$\frac{\sqrt{2}}{2}$};
    \draw (-0.06,-0.7071) -- (0.06,-0.7071) node[right=3pt] {\scriptsize$-\frac{\sqrt{2}}{2}$};
    \draw (-0.06,-1) -- (0.06,-1) node[left=3pt] {\scriptsize$-1$};
\end{tikzpicture}
\end{center}

\subsubsection*{Präsenzaufgabe 2: Allgemeines Verhalten von Schnitten unter Abbildungen}
\notizquelleDatei{handschrift/Uebungsblatt 3_261117_094748.PDF}{10--11}

\begin{satz}[Bild des Schnitts zweier Teilmengen]
Sei $f:X\to Y$ eine Abbildung und seien $A_1,A_2\subset X$ Teilmengen. Dann gilt:
\begin{enumerate}
    \item $f(A_1\cap A_2)\subset f(A_1)\cap f(A_2)$.
    \item Ist $f$ injektiv, so gilt Gleichheit: $f(A_1\cap A_2)=f(A_1)\cap f(A_2)$.
\end{enumerate}
\end{satz}

\begin{proof}
\begin{enumerate}
    \item Sei $y\in f(A_1\cap A_2)$ beliebig. Dann existiert ein $x\in A_1\cap A_2$ mit $f(x)=y$. Wegen $x\in A_1$ folgt $y=f(x)\in f(A_1)$, und wegen $x\in A_2$ folgt $y=f(x)\in f(A_2)$. Also gilt $y\in f(A_1)\cap f(A_2)$.

    \item Wegen Teil~1 genügt es zu zeigen, dass $f(A_1)\cap f(A_2)\subset f(A_1\cap A_2)$ gilt. Sei dazu $y\in f(A_1)\cap f(A_2)$ beliebig.
    \begin{itemize}
        \item Wegen $y\in f(A_1)$ gibt es ein $x_1\in A_1$ mit $f(x_1)=y$.
        \item Wegen $y\in f(A_2)$ gibt es ein $x_2\in A_2$ mit $f(x_2)=y$.
    \end{itemize}
    Aus $f(x_1)=y=f(x_2)$ folgt wegen der Injektivität von $f$, dass $x_1=x_2$. Setzen wir $x:=x_1=x_2$, so gilt $x\in A_1\cap A_2$ und daher $y=f(x)\in f(A_1\cap A_2)$.
\end{enumerate}
\end{proof}
